% Generated by scripts/build_cv_variants.py from cv-latex/variants.json.
% Edit the bilingual content source, then regenerate.
\documentclass[10pt,a4paper]{article}

\usepackage[a4paper,top=10mm,bottom=9mm,left=13mm,right=13mm,footskip=8mm]{geometry}
\usepackage{fontspec}

\usepackage{xcolor}
\usepackage[protrusion=true]{microtype}
\usepackage{ragged2e}
\usepackage{enumitem}
\usepackage{titlesec}
\usepackage{tabularx}
\usepackage{array}
\usepackage{needspace}
\usepackage{fancyhdr}
\usepackage{lastpage}
\usepackage{hyperref}

\definecolor{ink}{HTML}{111111}
\definecolor{inkii}{HTML}{3A3A3A}
\definecolor{inkiii}{HTML}{6A6A6A}
\definecolor{rule}{HTML}{8C8C8C}

% Fonts are found through OPENTYPEFONTS, which scripts/compile_cv.py points
% at cv-latex/fonts/.
\setmainfont{SourceSerif4SmText-Regular.otf}[
  BoldFont=SourceSerif4SmText-Bold.otf,
  Numbers=Lining,
  ItalicFont=SourceSerif4SmText-Regular.otf]
\newfontfamily\namefont{SourceSerif4Subhead-Bold.otf}[Numbers=Lining,BoldFont=SourceSerif4Subhead-Bold.otf]


\hypersetup{
  colorlinks=true,
  linkcolor=ink,
  urlcolor=ink,
  pdftitle={Lizhong Hu - Graduate Study CV},
  pdfauthor={Lizhong Hu},
  pdfsubject={Graduate Study CV},
  pdfcreator={XeLaTeX}
}

\renewcommand{\normalsize}{\fontsize{9.8}{11.8}\selectfont}
\renewcommand{\small}{\fontsize{8.62}{10.38}\selectfont}
\renewcommand{\footnotesize}{\fontsize{8.62}{10.38}\selectfont}
\normalsize
\AtBeginDocument{\color{ink}}

\setlength{\parindent}{0pt}
\setlength{\parskip}{0pt}
\setlist[itemize]{leftmargin=1.05em,labelsep=0.5em,itemsep=0.4pt,topsep=1pt,parsep=0pt,partopsep=0pt,
  label={\raisebox{0.2em}{\rule{0.28em}{0.28em}}}}

% Section headings: bold, spaced capitals (or spaced Chinese), a grey hairline.
\titleformat{\section}{\fontsize{9.8}{11.8}\selectfont\bfseries}{}{0pt}{\addfontfeatures{LetterSpace=8}\MakeUppercase}
  [{\vspace{1pt}\color{rule}\titlerule[0.4pt]}]
\titlespacing*{\section}{0pt}{4pt}{2pt}
\renewcommand{\arraystretch}{1.02}
\newcolumntype{Y}{>{\raggedright\arraybackslash}X}

% Quiet secondary text in place of italics (the site uses no italics).
\newcommand{\subtle}[1]{\leavevmode{\color{inkii}#1}}

% Entry heading: title in bold, date in grey on the right.
\newlength{\entrydatewidth}
\newcommand{\entry}[2]{%
  \par\addvspace{3pt}\noindent
  \if\relax\detokenize{#2}\relax
    \textbf{#1}\par
  \else
    \settowidth{\entrydatewidth}{{\small #2}}\addtolength{\entrydatewidth}{1.5pt}%
    \begin{tabularx}{\textwidth}{@{}Y@{\hspace{1.2em}}>{\raggedleft\arraybackslash}p{\entrydatewidth}@{}}
      \textbf{#1} & \leavevmode{\small\color{inkii}#2}
    \end{tabularx}\par
  \fi
  \nobreak\vspace{-0.5pt}
}

\newcommand{\skillrow}[2]{%
  \par\noindent\hangindent=1em\hangafter=1\textbf{#1: }#2\par
}

\pagestyle{empty}

\begin{document}
\RaggedRight
\setlength{\emergencystretch}{1em}
{\namefont\bfseries\fontsize{20}{22}\selectfont Lizhong Hu}\par
\vspace{2pt}{\small\color{inkii}\href{https://au1bhi.com}{au1bhi.com}\enspace\textperiodcentered\enspace \href{mailto:admin@au1bhi.com}{admin@au1bhi.com}\enspace\textperiodcentered\enspace \href{https://github.com/au1bhi}{GitHub}\enspace\textperiodcentered\enspace \href{https://codeforces.com/profile/Au1Bhi}{Codeforces}\enspace\textperiodcentered\enspace \href{https://www.linkedin.com/in/lizhong-hu-686b9b34a}{LinkedIn}}\par
\vspace{3pt}{\small\bfseries Information Retrieval · Verifiable RAG · Thesis Evaluation}\par
\vspace{3pt}{\color{ink}\rule{\textwidth}{0.7pt}}\par
\vspace{-8pt}
\section{Profile}
Software Engineering undergraduate in a Sino-American 4+0 dual-degree program (expected graduation: 2027). Interested in graduate study in computing, with a focus on information retrieval, verifiable RAG, and reliable software systems. Experience in thesis development, backend projects, and competitive programming.

\section{Education}
\begin{tabularx}{\textwidth}{@{}Y>{\raggedleft\arraybackslash}p{0.23\textwidth}@{}}
  \textbf{Quanzhou University of Information Engineering} & \subtle{2023--2027 (expected)} \\
  Software Engineering (Sino-foreign Cooperation) & GPA: 3.10/4.00 \\[0.5pt]
  \multicolumn{2}{@{}l@{}}{\textbf{Slippery Rock University}} \\
  Information Systems, BS (in progress; expected 2027) & GPA: 3.583/4.000
\end{tabularx}
\par\vspace{2pt}{\footnotesize\color{inkiii}Sino-American 4+0 dual-degree program conducted in China; approved by the Ministry of Education of China.}
\section{Research Interests}
Algorithms and data structures; software and distributed systems; information retrieval and retrieval-augmented generation; reliable intelligent information systems.
\section{Thesis \& Related Projects}
\entry{\href{https://github.com/au1bhi/NoteLLM}{NoteLLM: Personal Learning QA with Verifiable RAG}}{Ongoing}
\subtle{Undergraduate thesis}
\begin{itemize}
  \item Built a FastAPI backend with async SQLAlchemy, PostgreSQL, and pgvector for cosine-similarity retrieval; implemented page-aware extraction and overlapping chunks for PDF, Markdown, and TXT.
  \item Validated citations against retrieved chunks on the server and stored source excerpts for verification.
  \item A single baseline on 7 synthetic documents and 34 fixed questions recorded 100\% source-level Recall@5, 97.1\% automatic citation-to-source match, and 339 ms mean retrieval latency. These measure retrieval and citation-source matching on that synthetic set, not answer accuracy or general-domain performance.
  \item Public code and evaluation: github.com/au1bhi/NoteLLM; results in docs/evaluation/latest-results.md.
\end{itemize}
\entry{ETIPlus: Community Governance \& WeChat Management Backend}{Mar–Apr 2025}
\begin{itemize}
  \item Built Flask APIs with PostgreSQL and MongoDB for credit scoring, anomaly audits, and group management; ran Postman regression checks.
\end{itemize}
\section{Algorithms \& Community Contributions}
\entry{\href{https://www.matiji.net/exam/noi}{Matiji NOI-Pre Problem Bank | Volunteer Problem Tester}}{Jul–Oct 2024}
\begin{itemize}
  \item Created boundary test cases, validated solutions through targeted submissions, and reported defects in test data and judge results for the Matiji/CCF problem bank. \allowbreak\enspace{\small \mbox{\href{https://www.matiji.net/exam/volunteer}{Member showcase}}}
\end{itemize}
\entry{\href{https://new.saikr.com/vse/adpc/2025/spring}{6th College Algorithm Design \& Programming Challenge (Spring)}}{Spring 2025}
\subtle{One of the problem setters}
\begin{itemize}
  \item Contributed to problem design, solution validation, difficulty calibration, and editorial review.
\end{itemize}
\entry{University ICPC/CCPC Training Team | Coordinator}{Sep 2024–Sep 2026}
\begin{itemize}
  \item Coordinated university ICPC/CCPC team recruitment, training, selection, and participation logistics.
  \item Supported university teams at the 2026 ICPC Nanchang Invitational; 2025 CCPC Nanchang Invitational/Jiangxi Provincial Contest; 2026 CCPC Fujian Invitational/Provincial Contest; 2025 ICPC Asia EC and CCPC online contests; 2026 CCPC Guizhou Provincial Contest/National Invitational; 2025 ICPC Guizhou Provincial Contest; and 2024 ICPC Kunming Regional Contest.
\end{itemize}
\entry{Campus Programming Contest Organization \& Problem Setting}{}
\begin{itemize}
  \item ISCPC-AC Cup collegiate programming contest: principal organizer and problem setter.
  \item 2025 QUIE-Software School ACM Algorithm Contest (2nd Aipin Huiying Science \& Technology Festival): one of the organizers and problem setters.
\end{itemize}
\section{Competitive Programming \& Awards}
\begin{itemize}
  \item \href{https://au1bhi.com/files/WorldQuant-Challenge-Silver-Certificate.pdf}{WorldQuant Challenge (WorldQuant BRAIN): Silver Level, Nov 2026.}
  \item 2026 National Invitational of CCPC (Fujian), 13th Fujian Collegiate Programming Contest: First to Solve Award (Problem A; problem-level award).
  \item The 2024 ICPC Asia Kunming Regional Contest: Honorable Medal.
  \item 15th Lanqiao Cup: First Prize, Fujian Provincial Round, Python Programming (University Group B), 2024.
\end{itemize}
\section{Technical Skills}
\skillrow{Languages \& algorithms}{C++, Python, Java, JavaScript, Go; algorithms and data structures.}
\skillrow{Backend \& databases}{FastAPI, Flask, RuoYi, PostgreSQL, MongoDB, MySQL; REST API design and integration.}
\skillrow{Retrieval \& AI systems}{RAG, LLM API integration, prompt design.}
\skillrow{Systems \& operations}{Linux, Docker, Git, DOMjudge, Cloudflare Workers/CDN/DNS, WARP egress routing.}
\skillrow{Contest tools}{Familiar with Polygon (Codeforces) and Nowcoder; problem preparation, test-data management, and solution verification.}
\end{document}
